\chapterhead{4}{p-adic representations and congruences for modular forms}

\sechead{4.1}{p-adic representations and locally free sheaves}

Let $q$ be a power of $p$, $k$ a perfect field containing $\F_q$, $\W_n(k)$
its ring of Witt vectors\index{Witt vector} of length $n$, and $S_n$ a flat affine $\W_n(k)$-scheme
whose special fibre is normal, reduced and irreducible. Suppose that $S_n$
admits an endomorphism $\varphi$ which induces the q-th power mapping on the
special fibre. [If $S_n$ is affine and smooth over $\W_n(k)$, then such a
$\varphi$ always exists.]

\begin{kathm}{Proposition}{4.1.1}
There is an equivalence of categories between the category of finite free
$\W_n(\F_q)$-modules $M$ on which $\pi_1(S_n)$ acts continuously, and the
category of pairs $(H,F)$ consisting of a locally free sheaf of finite rank
$H$ on $S_n$ together with an isomorphism $F\colon\varphi^*(H)\to H$.
\end{kathm}

\begin{kaproof}[Construction-proof]
Given a representation $M$ of $\pi_1(S_n)$, let $T_n$ be a finite \'etale galois
$S_n$-scheme such that the representation factors through $\Aut(T_n/S_n)$.
Because $T_n$ is \'etale over $S_n$, there is a unique $\varphi$-linear
endomorphism of $T_n$ which induces the q-th power endomorphism of
$T_n\times_{\W_n(k)}k$, which we denote by $\varphi_T$. By unicity, $\varphi_T$
commutes with $\Aut(T_n/S_n)$. Let $H_T$ be the $T_n$-module
$M\tens_{\W_n(\F_q)}\cO_{T_n}$, and let $F_T$ be the $\varphi_T$-linear
endomorphism of $H_T$ defined by $F_T(m\tens f)=m\tens\varphi_T(f)$. For each
$g\in\Aut(T_n/S_n)$, we define $g(m\tens f)=g(m)\tens(g^{-1})^*(f)$, thus
defining an action of $\Aut(T_n/S_n)$ on $(H_T,F_T)$. By descent, it follows
that there is a unique $(H,F)$ on $S_n$ whose inverse image on $T_n$ is
$\Aut(T_n/S_n)$-isomorphic to $(H_T,F_T)$. The construction
$M\rightsquigarrow(H,F)$ defines the functor we will prove to be an equivalence.
Notice that we can recover $M$ as the fixed points of $F_T$ acting as
$\varphi$-linear endomorphisms of the module of global sections of $H_T$, hence
our functor is fully faithful. To show that it is an equivalence, we must show
that any $(H,F)$ arises in this way, or, in concrete terms, we must show that
given $(H,F)$, there exists a finite \'etale covering $T_n$ of $S_n$ over which
$H$ admits a basis of F-fixed points. We proceed by induction on the integer
$n$.

Suppose first $n=1$. Then $S$ is a k-scheme, and $(H,F)$ is a locally free
finite rank S-module $H$ together with a q-linear endomorphism $F$ of $H$ which
gives an isomorphism $F\colon H^{(q)}\to H$. For any S-scheme $T$, the inverse
image module $H_T$ carries the inverse image q-linear map $F_T$, defined by
$(F_T(h\tens t)=F(h)\tens t^q$, which gives an isomorphism
$F_T\colon H_T^{(q)}\to H_T$.

Notice that the functors on S-schemes
\[
\begin{cases}
X(T) = \text{global sections of } H_T \\
Y(T) = \text{bases of } H_T\ (\cO_T\text{-isomorphisms }(\cO_T)^r\to H_T,\text{ where } r=\rank(H)) \\
Z(T) = \text{bases of } H_T \text{ consisting of fixed points of } F_T
\end{cases}
\]
are all \uline{representable}, the first by $\Spec_S(\Sym(H))$, the second by
the open subset of the $r=\rank(H)$-fold product
$X^{(r/S)}=X\times_S X\times\cdots\times_S X$ over which the tautological map
$(\cO_{X^{(r/S)}})^r\to H_{X^{(r/S)}}$ is an isomorphism, the third by the
closed subscheme of $Y$ over which the universal basis is fixed by $F_Y$. We
must show that $Z$ is finite and \'etale over $S$. \uline{This} problem is local
on $S$, hence we may assume $S$ affine and $H$ free. Choose a basis
$h_1,\dots,h_r$ of $H$, and let $(a_{ij})$ be the invertible matrix of
$F\colon F(h_i)=\sum_j a_{ji}h_j$.

Consider the functor on S-schemes
\[
Y'(T) = \text{sections of } H_T \text{ fixed by } F_T.
\]
It is representable by a scheme finite and \'etale of rank $q^r$ over $S$,
because a section $\sum_i X_ih_i$ of $H$ is F-fixed if and only if
$\sum_j X_jh_j=\sum_i (X_i)^q\sum_j a_{ji}h_j$, thus $Y'$ is the closed
subscheme of $\mathbb{A}^r_S$ defined by the equation
\[
X_j = \sum_i a_{ji}(X_i)^q,\qquad j=1,\dots,r .
\]
Because the matrix $(a_{ij})$ is invertible, if we denote by $(b_{ij})$ its
\uline{inverse}, the equations are the same as the equations
\[
(X_i)^q = \sum_j b_{ij}X_j\qquad i=1,\dots,r ,
\]
which define a finite \'etale S-scheme of rank $q^r$.

The scheme $Z$ is the \uline{open} subscheme of $Y'^{(r/S)}=Y'\times_S\cdots
\times_S Y'$ where the universal r-tuple of F-fixed sections form a base of
$H$, and hence $Z$ is \uline{\'etale} over $S$. It remains to check that $Z$ is
proper over $S$, and non-void. By the valuative criterion, we must show that for
any valuation ring $V$ over $S$, any F-fixed basis of $H_K$ ($K$ the fraction
field of $V$) prolongs to an F-fixed basis of $H_V$. Because the scheme $Y'$ of
fixed points is \uline{finite} over $S$, each basis element prolongs to a
unique F-fixed section of $H_V$. To see that the corresponding map
$V^r\to H_V$ is an isomorphism, we look at its \uline{determinant}, which
reduces us to the case of a rank one module. Then the matrix of $F$ is
$F(h_1)=ah_1$, with $a$ invertible in $V$, and an F-fixed basis of $H_K$ is a
vector $k\cdot h_1$, with $k\in K$ satisfying $k=ak^p$. As $a\in V$ is
invertible in $V$, any such $k$ is an invertible element of $V$, hence
$k\cdot h_1$ "is" an F-fixed base of $H_V$.

It remains to see that $Z$ is non-empty. As its formation commutes with
arbitrary change of base $S'\to S$, it's enough to check the case when $S$ is
the spectrum of an algebraically closed field. But a finite-dimensional vector
space over an algebraically closed field with a q-linear \uline{automorphism}
is always spanned by its fixed points (Lang's trick; cf.\cite{katz69}) and the set of
fixed bases is a $\GL_r(\F_q)$-torseur. Thus $Z$ is finite \'etale of rank
$=\#\GL_r(\F_q)$ over $S$, and the action of $\GL_r(\F_q)$ on $Z$ (induced by
its action on the functor of F-fixed bases) makes $Z$ into a
$\GL_r(\F_q)_S$-torseur. The cohomology class of this torseur is an element of
$H^1_{\text{\'et}}(S,\GL_r(\F_q))=\Hom(\pi_1(S),\GL_r(\F_q))$ which is none
other than the desired representation. This concludes the construction-proof
for $n=1$.

Suppose the result known for $n-1$. Then there is a finite \'etale covering
$T_{n-1}$ of $S_{n-1}=S_n\times_{\W_n(k)}\W_{n-1}(k)$ over which $H/p^{n-1}H$
admits a basis of F-fixed points. There is a unique finite \'etale covering
$T_n$ of $S_n$ such that $T_n\times_{S_n}S_{n-1}$ is $T_{n-1}$, and replacing
$S_n$ by $T_n$ we may suppose that $H/p^{n-1}H$ admits a basis of F-fixed
points. Let $h_1,\dots,h_r$ be a basis of $H$ which \uline{lifts} an F-fixed
basis of $H/p^{n-1}H$ ($S_n$ is affine!). Writing
$\underline{h}=\begin{pmatrix}h_1\\ \vdots\\ h_r\end{pmatrix}$, we have
$F(\underline{h})=(1+p^{n-1}\Delta)\underline{h}$. In order for
$(1+p^{n-1}E)\cdot\underline{h}$ to be an F-fixed basis, we must have
\[
(1+p^{n-1}\cdot\varphi(E))\cdot(1+p^{n-1}\Delta)\underline{h} = (1+p^{n-1}E)\underline{h}
\]
or equivalently ($S_n$ being flat over $\W_n(k)$)
\[
\varphi(E)+\Delta\equiv E \bmod (p),
\]
which is a set of $r^2$ Artin-Schreier equations $(e_{ij})^q-e_{ij}=-\Delta_{ij}$
over $S_1=S_n\times_{\W_n(k)}k$. On a finite \'etale covering $T_1$ of $S_1$,
these equations admit solutions, and hence on the unique finite \'etale covering
$T_n$ of $S_n$ such that $T_n\times_{S_n}S_1=T_1$, the module $H_{T_n}$ admits
an F-fixed basis.
\end{kaproof}

\begin{kathmstar}[4.1.2.1]{Remarks 4.1.2.1}
The operation "tensor product" in the category of representations of
$\pi_1(S_n)$ in finite free $\W_n(\F_q)$ modules corresponds to the tensor
product $(H,F)\tens(H',F')=(H\tens_{S_n}H',F\tens F')$, defined by
$(F\tens F')(h\tens h')=F(h)\tens F'(h')$.
\end{kathmstar}

\kaitem{4.1.2.2} The "internal Hom" in the category of representations
corresponds to the internal Hom defined by
$\underline{\Hom}((H,F),(H_1,F_1))=(\underline{\Hom}(H,H_1),F_2)$ where $F_2$ is
the unique $\varphi$-linear endomorphism of $\underline{\Hom}(H,H_1)$ such that
for $h\in H$, $f\in\underline{\Hom}(H,H_1)$, we have $F_2(f)(F(h))=F_1(f(h))$.
In particular, $\underline{\Hom}((H,F),(\cO,\varphi))$ is the "contragredient"
$(\widetilde{H},\widetilde{F})$, defined by the requirement that for $h\in H$,
$\widetilde{h}\in\widetilde{H}$, we have
$\langle F(h),\widetilde{F}(\widetilde{h})\rangle=\varphi(\langle h,\widetilde{h}\rangle)$.

\kaitem{4.1.2.3} Because $S_1$ is normal, reduced and irreducible, a
representation of $\pi_1(S_m)=\pi_1(S_1)$ is just a suitably unramified
representation of the Galois group of the function field of $S_1$. Thus for any
non-void open set $U\subset S_n$, the functor "restriction" from the category of
representations of $\pi_1(S_n)$ to the category of representations of
$\pi_1(U)$ is \uline{fully faithful}. Hence the functor "restriction" from the
category of $(H,F)$'s over $S_n$ to the category of those over $U$ is \uline{fully
faithful}.

\sechead{4.2}{Application to modular schemes}

\kaitemb{4.2.0.} Let $n\geq 3$, $p$ a prime not dividing $n$, $q$ a power of $p$
such that $q\equiv 1 \bmod n$, and choose an \uline{isomorphism} between $\mu_n$
and $\Z/n\Z$ over $\W(\F_q)$, i.e. choose a primitive n'th root of unity
$\zeta$. Let $S_m^\zeta$ (resp. $\overline{S}_m^\zeta$) be the open subset of
$M_n\tens\W_m(\F_q)$ (resp. of $\Mbar_n\tens\W_m(\F_q)$) where $E_{p-1}$ is
invertible \uline{and} where the e.m. pairing on the basis of ${}_nE$ has the
value $\zeta$, i.e. where the \uline{determinant} of the level n structure is
the chosen isomorphism of $\Z/n\Z$ with $\mu_n$. The schemes $S_m^\zeta$ (resp.
$\overline{S}_m^\zeta$) are smooth affine $\W_m(\F_q)$ schemes with
geometrically connected fibres. In the notation of \eqref{ka:2.9}, we have
$M_n(\W_m(\F_q),1)=\bigcup S_m^\zeta$, the union taken over the \uline{primitive}
n'th roots of unity, and $\Mbar_n(\W_m(\F_q),1)=\bigcup\overline{S}_m^\zeta$.

Let $\sigma$ denote the Frobenius\index{Frobenius} automorphism of $\W_m(\F_q)$. We have
$\sigma(\zeta)=\zeta^p$, and hence $S_m^{\zeta^p}=(S_m^\zeta)^{(\sigma)}$,
$\overline{S}_m^{\zeta^p}=(\overline{S}_m^\zeta)^{(\sigma)}$. The endomorphism
$\varphi$ of $\Mbar_n(\W_m(\F_q),1)$ defined by "division by the canonical
subgroup" does not respect the various $\overline{S}_m^\zeta$, but rather it
maps $\overline{S}_m^\zeta$ to $\overline{S}_m^{\zeta^p}$ (because modulo $p$,
the canonical subgroup\index{canonical subgroup} is the kernel of absolute Frobenius). As
$\overline{S}_m^{\zeta^p}=(\overline{S}_m^\zeta)^{(\sigma)}$, we may and will
view $\varphi$ as a $\sigma$-linear endomorphism of each $S_m^\zeta$, which
modulo $p$ becomes the p'th power mapping. In a similar fashion, the
endomorphism $\varphi$ of the invertible sheaf $\w^{\otimes k}$ on
$\Mbar_n(\W_m(\F_q),1)$, defined by
$\varphi(f)(E,\omega,\alpha_n)=f(E/H,\pi^*(\omega),\pi(\alpha_n))$ [where $H$
denotes the canonical subgroup and $\pi\colon E\to E/H$ the projection], may be
viewed as a $\varphi$-linear endomorphism of $\w^{\otimes k}|S_m^\zeta$, for
each primitive n'th root of unity $\zeta$. [Notice that $\w^{\otimes k}$ is
generated by $\varphi(\w^{\otimes k})$ as a sheaf; indeed for a local section
$f$ of $\w^{\otimes k}$, a glance at q-expansions shows that
$\varphi(f)\equiv f^p(E_{p-1})^k$, hence $\varphi(f)$ is an invertible section
wherever $f$ is.]

We wish to determine \uline{which} representation of $\pi_1(\overline{S}_m^\zeta)$
in a free $\Z/p^m\Z=\W_m(\F_p)$-module of rank one corresponds via \eqref{ka:4.1.1} to
$(\w^{\otimes k},\varphi)$ on $\overline{S}_m^\zeta$. Of course it suffices to
do this for $k=1$, by \eqref{ka:4.1.2.1}. There is an obvious candidate, namely the
representation of $\pi_1(\overline{S}_m^\zeta)$ on the \uline{\'etale quotient}
of the kernel of $p^m$ on the \uline{universal} curve $E$. [Noting by
$\pi\colon E\to E^{(\varphi)}=E/H$ the projection onto the quotient by the
canonical subgroup, the composite $\pi_m\colon E\to E^{(\varphi^m)}$ induces an
isomorphism of the \'etale quotient
${}_{p^m}E/{}_{p^m}\widehat{E}={}_pE/\Ker(\pi^m)\xrightarrow{\sim}
\Ker(\widetilde{\pi})^m$ in $E^{(\varphi^m)}$.] If this candidate is to "work",
we must have:

\begin{kathm}{Lemma}{4.2.1.}
The representation of $\pi_1(S_m^\zeta)$ on $\Ker(\widetilde{\pi})^m$ extends to
a representation of $\pi_1(\overline{S}_m^\zeta)$, i.e., it is "unramified at
$\infty$".
\end{kathm}

\begin{kaproof}
Since the \'etale topology cannot distinguish $\overline{S}_m^\zeta$ and
$\overline{S}_1^\zeta$, it is equivalent to show that the representation of
$\pi_1(S_1^\zeta)$ on $\Ker(V^m)$ extends to a representation of
$\pi_1(\overline{S}_1^\zeta)$ on $\Ker(V^m)$. Let $K$ denote the function field
of $\overline{S}_1^\zeta$; we must see that the inertia group of
$\Gal(K^{\mathrm{sep}}/K)$ at each cusp acts trivially on $\Ker(V^m)$ in
$E_K^{(\varphi^m)}(K^{\mathrm{sep}})$. To decide, we may replace $K$ by its
\uline{completion} at each cusp, which is just $k((q))$, $k=\F_q$, and the
inverse image of $E$ over this completion is the Tate curve\index{Tate curve}
$\mathrm{Tate}(q^n)/k((q))$. The curve $E^{(\varphi^m)}$ becomes
$\mathrm{Tate}(q^{np^m})$, and $(\widetilde{\pi})^m$ is the map
$\mathrm{Tate}(q^{np^m})\to\mathrm{Tate}(q^n)$ given by "division by the
subgroup generated by $q^n$". As this subgroup consists entirely of rational
points, the inertial group (and even the decomposition group) at each cusp acts
trivially.
\end{kaproof}

\begin{kathm}{Theorem}{4.2.2.}
The representation of $\pi_1(S_m^\zeta)$ on $\Ker(\widetilde{\pi})^m$
($\simeq$ to the \'etale quotient of $\Ker p^m$ on the universal curve)
corresponds, via the equivalence \eqref{ka:4.1.1}, to $(\w,\varphi)$.
\end{kathm}

\begin{kaproof}
By the "full-faithfulness" of restriction to open sets, it suffices to prove
this over $\overline{S}_m^\zeta$. Let's take a finite \'etale covering $T$ of
$S_m^\zeta$ which trivializes the representation - in concrete terms, we adjoin
the coordinates of a point of $\Ker(\widetilde{\pi})^m$ of order precisely
$p^m$. Over $T$, each point of $\Ker(\widetilde{\pi})^m$ gives a morphism
$(\Z/p^m\Z)_T\xrightarrow{\sim}(\Ker(\widetilde{\pi})^m)_T$, whose \uline{Cartier
dual} is a morphism $(\Ker p^m \text{ in }\widehat{E})_T=(\Ker\widetilde{\pi}^m)_T
\to(\mu_{p^m})_T\xrightarrow{\sim}(\Gm)_T$. The inverse image of the invariant
differential $dt/t$ on $(\Gm)_T$ furnishes an invariant differential on the
kernel of $p^m$ in $\widehat{E}$. Since $T$ is killed by $p^m$, the first
infinitesimal neighborhood of the identity section of $E$ lies in the kernel of
$p^m$ in $\widehat{E}$, and hence there is a unique invariant differential on
$E$ whose restriction to the kernel of $p^m$ in $\widehat{E}$ is the given one.
Thus we have defined a \uline{morphism} from $(\Ker(\widetilde{\pi})^m)_T$ to
$\w_T$. Further, if we take a point of $\Ker(\widetilde{\pi})^m$ of order
\uline{precisely} $p^m$, the map $(\Z/p^m\Z)_T\to(\Ker(\widetilde{\pi})^m)_T$
is an \uline{isomorphism}, hence the Cartier dual is an isomorphism, and hence
the inverse image of $dt/t$ on $\Ker p^m$ in $\widehat{E}$ is nowhere vanishing.
Thus the induced map
$(\Ker(\widetilde{\pi})^m)_T\tens_{\Z/p^m\Z}\cO_T\to\w_T$ is an isomorphism of
invertible sheaves on $T$. It is clear that this map commutes with the obvious
action of $\Aut(T/S_m^\zeta)$. [In concrete terms, and locally on $S$,
$\Ker(p^m)$ in $\widehat{E}$ has coordinate ring free on $1,X,\dots,X^{p^m-1}$,
a point $P$ of $(\Ker(\widetilde{\pi})^m)_T$ gives rise to a map $\mu_{p^m}$
defined by $f(X)=\sum a_i(P)X^i$, the corresponding differential is
$\omega_P=df/f$, and for any $g\in\Aut(T/S_m)$, we have $a_i(g(P))=g(a_i(P))$,
and hence $\omega_{g(P)}=g(\omega_P)$.] By descent, we have constructed an
isomorphism between $\w$ and the invertible sheaf on $S_m^\zeta$ associated to
the \'etale quotient of ${}_{p^m}E$.

It remains to see that this isomorphism is compatible with the $\varphi$-linear
endomorphisms. Tensoring one with the inverse of the other, we obtain a
$\varphi$-linear endomorphism on $\cO_{S_m}$; we must show that it carries "1"
to "1".

To check this, it suffices to do so in a "punctured disc at $\infty$", over
$\W_m(\F_q)((q))$ when we look at the Tate curve $\mathrm{Tate}(q^n)$. The
morphism $\widetilde{\pi}\colon\mathrm{Tate}(q^{n\cdot p^m})\to\mathrm{Tate}(q^n)$
has kernel the subgroup generated by $q^n$. The point $q^n$ is a rational point
of $\Ker(\widetilde{\pi})^m$, and the corresponding differential is precisely
the Tate differential $\omega_{\mathrm{can}}=dt/t$. As $q^n$ is a rational
point, the section $[q^n]\tens 1$ of
$\Ker(\widetilde{\pi})^m\tens_{\Z/p^m\Z}\cO$ is fixed by the canonical $F$, and
the corresponding section $\omega_{\mathrm{can}}$ of $\w$ is fixed by $\varphi$
(because $\omega_{\mathrm{can}}$ has q-expansion identically "1"). Hence our
isomorphism respects the $\varphi$-linear endomorphisms in a punctured disc
around $\infty$, and hence respects it everywhere.
\end{kaproof}

\begin{kathmstar}{Remark 4.2.2.1.}
One may prove this theorem in a non-constructive way by showing that both of the
associated p-adic characters $\chi_i\colon\pi_1(S_m^\zeta)\to(\Z/p^m\Z)^\times$
have the \uline{same value} on all \uline{Frobenius} elements, namely the
reciprocal of the "unit root" of the ordinary elliptic curve\index{elliptic curve} which is the fibre
over the corresponding closed point of $S_m^\zeta$.
\end{kathmstar}

\secmark{4.3}{Igusa's theorem}

\begin{kathm}{Theorem}{4.3.}
(Igusa \cite{igusa68}) The homomorphism
$\pi_1(\overline{S}_m^\zeta)\to\Aut(\Ker(\widetilde{\pi})^m)\xrightarrow{\sim}
(\Z/p^m\Z)^\times$ is \uline{surjective}, and for every non-void open set
$U\subset\overline{S}_m^\zeta$, the composite
$\pi_1(U)\to\pi_1(\overline{S}_m^\zeta)\to(\Z/p^m\Z)^\times$ remains
\uline{surjective}.
\end{kathm}

\begin{kaproof}
It suffices to show that, denoting by $K$ the function field of
$S_m^\zeta\times_{\W_m(\F_q)}\F_q$, the homomorphism
$\Gal(K^{\mathrm{sep}}/K)\to\Aut(\Ker V^m \text{ in }E^{(p^m)}(K^{\mathrm{sep}}))$
is surjective. In fact, we will prove that the inertial group of
$\Gal(K^{\mathrm{sep}}/K)$ at any \uline{supersingular} elliptic curve already
maps surjectively. Let $\varphi$ be any closed point of $S_1^\zeta$ where
$E_{p-1}$ vanishes; replacing $\F_q$ by its algebraic closure $k$, we may assume
$\varphi$ is a rational point. The completion of $S_1^\zeta\tens k$ at $\varphi$
is isomorphic to $\Spec(k[[A]])$, and the inverse image of the universal curve
over $k[[A]]$ admits a nowhere vanishing differential $\omega$ such that
$E_{p-1}(E,\omega)=A$. (This is just Igusa's theorem\index{Igusa's theorem} that the Hasse invariant\index{Hasse invariant}
has only simple zeros.) So we must prove
\end{kaproof}

\begin{kathm}{Theorem}{4.3 bis}
(Igusa) Let $E,\omega$ be an elliptic curve over $k[[A]]$ with Hasse invariant
$A$, $k$ being an algebraically closed field of characteristic $p$. Then the
extension of $k((A))$ obtained by adjoining the points of
$\Ker V^m\colon E^{(p^m)}\to E$ is fully \uline{ramified} of degree
$p^{m-1}(p-1)$, with Galois group canonically isomorphic to
$\Aut(\Z/p^m\Z)$.
\end{kathm}

\begin{kaproof}
The first statement implies the second, since $\Ker V^m$ is cyclic of order
$p^m$ over $k((A))^{\mathrm{sep}}$. In terms of a normalized parameter $X$ for
the formal group\index{formal group} (i.e. $[\zeta](X)=\zeta X$ for any p-1'st root of unity
$\zeta\in\Z_p$), the endomorphism $[p]$ has the shape
\begin{kadispb}{4.3.1}
[p](X) = V(X^p) = AX^p + \alpha X^{p^2} + \cdots
\end{kadispb}
with $\alpha$ invertible in $k[[A]]$ (because modulo $A$, we have a
supersingular curve by hypothesis, hence its formal group is of height
\uline{two}). Thus $V(X)=AX+\alpha X^p+\cdots$, and the composite
$V^m\colon E^{(p^m)}\to E$ is the composite
\[
E^{(p^m)}\xrightarrow{V^{(p^{m-1})}}E^{(p^{m-1})}\xrightarrow{V^{(p^{m-2})}}
\cdots\xrightarrow{V^{(p)}}E^{(p)}\xrightarrow{V}E .
\]
The expression of $V^{(p^i)}$ is $V^{(p^i)}(X)=A^{p^i}X+\alpha^{p^i}X^p+\cdots$.
A point of $\Ker V^m$ with values in $k((A))^{\mathrm{sep}}$ of order precisely
$p^m$ may be viewed as a sequence $y_0,\dots,y_{m-1}$ of elements of the maximal
ideal of $k((A))^{\mathrm{sep}}$ which satisfy the successive equations
\[
\begin{cases}
0 = V(y_0) = Ay_0 + \alpha(y_0)^p + \cdots \\
y_0 = V^{(p)}(y_1) = A^py_1 + \alpha^p(y_1)^p + \cdots \\
y_{m-2} = V^{(p^{m-1})}(y_{m-1}) = A^{p^{m-1}}y_{m-1} + \alpha^{p^{m-1}}(y_{m-1})^p + \cdots .
\end{cases}
\]
But a glance at the Newton polygons\index{Newton polygon} of these equations shows successively that
the \uline{ordinals} of $y_0,\dots,y_{m-1}$ are given by (noting by $\ord$ the
ordinal normalized so that $\ord(A)=1$):
\[
\begin{cases}
\ord(y_0) = 1/p-1 \\
\ord(y_1) = 1/p(p-1) \\
\qquad\vdots \\
\ord(y_{m-1}) = 1/p^{m-1}(p-1) .
\end{cases}
\]
\end{kaproof}

\sechead{4.4}{Applications to congruences between modular forms \`a la Serre}

\begin{kathm}{Corollary}{4.4.1.}
Let $k$ be an integer, and suppose $m\geq 1$. The following conditions are
equivalent:
\begin{enumerate}
\item[1)] $k\equiv 0 \bmod (p-1)\cdot p^{m-1}$ if $p\neq 2$, and
$k\equiv 0 \bmod 2^{\alpha(m)}$ if $p=2$, where $\alpha(1)=0$, $\alpha(2)=1$, and
$\alpha(m)=m-2$ if $m\geq 3$.
\item[2)] The k'th (tensor) power of the representation of
$\pi_1(\overline{S}_m^\zeta)$ on the \'etale quotient of ${}_{p^m}E$ is trivial.
\item[3)] The sheaf $\w^{\otimes k}$ on $\overline{S}_m^\zeta$ admits a nowhere
vanishing section fixed by $\varphi$.
\item[4)] Over a non-void open set $U\subset\overline{S}_m^\zeta$,
$\w^{\otimes k}$ admits a nowhere vanishing section fixed by $\varphi$.
\item[5)] Over $\overline{S}_m^\zeta$, $\w^{\otimes k}$ admits a section whose
q-expansion at one of the cusps of $\overline{S}_m^\zeta$ is identically 1.
\item[6)] Over a non-void open set $U\subset\overline{S}_m^\zeta$ which contains
a cusp, $\w^{\otimes k}$ admits a section whose q-expansion at that cusp is
identically 1.
\end{enumerate}
Further, if 1) holds, then any section verifying either 4) or 6) extends uniquely
to a section over all of $\overline{S}_m^\zeta$ verifying 3) and 5), and is in
fact the k/p-1'st power of $E_{p-1}$.
\end{kathm}

\begin{kaproof}
1) $\iff$ 2), because the image of $\pi_1(\overline{S}_m^\zeta)$ is all of
$\Aut(\Z/p^m\Z)\simeq(\Z/p^m\Z)^*$, a group of exponent $p^{m-1}(p-1)$ for
$p\neq 2$ and of exponent $2^{\alpha(m)}$ for $p=2$. By \eqref{ka:4.3}, 2) $\iff$ 3)
equivalence 3) $\iff$ 4) is by full-faithfulness of "restriction to U",
cf.\eqref{ka:4.1.2.3}. By the explicit formula for $\varphi$ and the q-expansion
principle, we have 3) $\iff$ 5) and 4) $\iff$ 6). When 1) holds, the unicity of
the section satisfying 4) or 6) or 3) or 5) follows from the full-faithfulness of
restriction to $U$; that this section is $(E_{p-1})^{k/p-1}$ follows from the
q-expansion principle\index{q-expansion principle}.
\end{kaproof}

\begin{kathm}{Corollary}{4.4.2.}
(Serre) Suppose $f_i$, $i=1,2$ are elements of $S(\W(\F_q),1,n,k_i)$, $i=1,2$,
and that $k_1\geq k_2$. Suppose that the q-expansions of $f_1$ and $f_2$ at at
least one cusp of $\Mbar_n(\W(\F_q),1)$ are congruent modulo $p^m$, and that
$f_1(q)\not\equiv 0 \bmod p$ at that cusp. Then $k_1\equiv k_2 \bmod
p^{m-1}(p-1)$ if $p\neq 2$, and $k_1\equiv k_2 \bmod 2^{\alpha(m)}$ if $p=2$,
where $\alpha(1)=0$, $\alpha(2)=1$, and $\alpha(m)=m-2$ for $m\geq 3$. If these
congruences hold at at least one cusp on each irreducible component, then
$f_2\equiv f_1\cdot(E_{p-1})^{k_2-k_1/p-1} \bmod p^mS(\W(\F_q),1,n,k_2)$.
\end{kathm}

\begin{kaproof}
Once we prove the congruence on the $k_i$, the final assertion results from the
q-expansion principle. To prove the congruence on the weights, we reduce the
situation modulo $p^m$. Then $f_1$ and $f_2$ are sections of $\w^{\otimes k_1}$
and $\w^{\otimes k_2}$ respectively over $\overline{S}_m^\zeta$. By hypothesis,
$f_1$ and hence $f_2$ are invertible on a non-void open set $U$ of
$\overline{S}_m^\zeta$, and the ratio $f_2/f_1$ is thus an invertible section of
$\w^{k_2-k_1}$ over $U$, and by hypothesis $f_2/f_1$ has q-expansion identically
one at at least one cusp on each $\overline{S}_m^\zeta$. By \eqref{ka:4.4.1}, we have the
desired congruence on $k_1-k_2$.
\end{kaproof}

\begin{kathm}{Corollary}{4.4.3.}
(Serre) Let $f$ be a true modular form\index{modular form} of level $n$ and weight $k$ on
$\Gamma_0(p)$, holomorphic at the unramified cusps, and defined over the
fraction field $K$ of $\W(\F_q)$. Suppose that at each unramified cusp, the
q-expansion has all its \uline{non-constant} q-coefficients in $\W(\F_q)$. Then
the \uline{constant terms} of the q-expansions lie in $p^m\cdot\W(\F_q)$, where,
for $p\neq 2$, $m$ is the \uline{largest} integer such that
$\varphi(p^m)=\#(\Z/p^m\Z)^\times$ divides $k$, and for $p=2$, $m=1$ if $k$ is
odd, and $m=\ord_2(k)+2$ if $k$ is even.
\end{kathm}

\begin{kaproof}
For $N\gg 0$, $p^Nf$ is a true modular form of level $n$ and weight $k$ on
$\Gamma_0(p)$, defined over $\W(\F_q)$. By \eqref{ka:3.2}, there is a unique element $g$
of $S(\W(\F_q),1,n,k)$ whose q-expansions are those of $p^Nf$ at the
corresponding unramified cusps. If $-m_0$ denotes the minimum of the ordinals of
the constant terms of these q-expansions, then $g$ is divisible by $p^{N-m_0}$ in
$S(\W(\F_q),1,n,k)$, by \eqref{ka:2.7}. Thus we may write $g=p^{N-m_0}h$ with
$h\in S(\W(\F_q),1,n,k)$ having the same q-expansions as does $p^{m_0}f$ at the
corresponding unramified cusps. Then $h$ has integral q-expansions, and at least
one of them is congruent modulo $p^{m_0}$ to a \uline{constant} which is a
\uline{unit} on $\W(\F_q)$. Multiplying $f$ by the reciprocal of this unit, we
get a q-expansion which is congruent mod $p^N$ to "1". As the constant function
"1" is modular of weight zero, we must have $k\equiv 0 \bmod p^{m_0-1}(p-1)$ if
$p\neq 2$, $k\equiv 0 \bmod 2^{\alpha(m_0)}$ for $p=2$.
\end{kaproof}

\begin{kathmstar}{Remark 4.4.4.}
If we apply these estimates to the constant terms of the classical level one
Eisenstein series $E_k$, we get \uline{precisely} the correct bounds for the
denominators of the Bernoulli numbers\index{Bernoulli number} (cf. \cite{serre-proceedings}, \cite{swinnerton-dyer-proceedings} for more on Bernoulli
numbers).
\end{kathmstar}

\sechead{4.5}{Applications to Serre's "modular forms of weight $X$"}

\kaitemb{4.5.0.} Let $X\in\End(\Z_p^\times)$. For each power $p^m$ of $p$, $X$
induces an endomorphism of $(\Z/p^m\Z)^\times$. For any primitive n'th root of
unity $\zeta$, and for any representation $\rho$ of $\pi_1(S_m^\zeta)$ in a free
$\Z/p^m\Z$ module of rank one, we may define the representation
$\rho^X\xlongequal{\text{defn}}X\circ\rho$. Taking for $\rho$ the representation
given by the \'etale quotient of ${}_{p^m}E$, we denote by $(\w^X,\varphi)$ the
invertible sheaf with $\varphi$-linear endomorphism which corresponds to
$\rho^X$. For variable $m$, the sheaves $\w^X$ on $S_m^\zeta$ are compatible,
and we define a compatible family of global sections to be a p-adic modular form\index{p-adic modular form}
of weight $X$ and level $n$, holomorphic at $\infty$, defined over $\W(\F_q)$.
If $X=k\in\Z\subset\End(\Z_p^\times)$, we just recover the elements of
$S(\W(\F_q),1,n,k)$. For $p\neq 2$, $\Z$ is dense in $\End(\Z_p^\times)$, and
indeed $\End(\Z_p^\times)\xleftarrow{\sim}\varprojlim\Z/\varphi(p^m)\Z$; for
$p=2$, $\Z_2$ has \uline{index four} in the (non-commutative) ring
$\End(\Z_2^\times)$. If $p\neq 2$, then for any $X$ (resp. if $p=2$, for any
$X\in\Z_2$), the pair $(\w^X,\varphi)$ on $\overline{S}_m^\zeta$ is isomorphic
to $(\w^{\otimes k_m},\varphi)$ for any $k_m\in\Z$ such that
$k_m\equiv X \bmod \varphi(p^m)$ (resp. if $p=2$, modulo $2^0$ if $m=1$, $2^1$
if $m=2$, and $2^{m-2}$ if $m\geq 3$). The isomorphism between
$(\w^{\otimes k_m},\varphi)$ and $(\w^{\otimes k_m'},\varphi)$ for different
choices $k_m,k_m'\in\Z$ approximating $X$ is given by multiplication by
$(E_{p-1})^{(k_m'-k_m)/(p-1)}$. As this isomorphism leaves \uline{invariant} the
q-expansion modulo $p^m$ (resp. modulo $2^{m-1}$ for $p=2$), it follows that a
p-adic modular form of weight $X$ and level $n$, holomorphic at $\infty$, defined
over $\W(\F_q)$, has a well defined q-expansion in $\W[[q]]$ at each cusp, and
that for given $X$, $f$ is uniquely determined by its q-expansions.

\begin{kathm}{Theorem}{4.5.1.}
Let $X\in\End(\Z_p^\times)$, and suppose $X\in\Z_2$ if $p=2$. Let $f$ be a
modular form\index{modular form} of weight $X$ and level $n$, holomorphic at $\infty$, defined over
$\W(\F_q)$. Then there exists a sequence of integers $0\leq k_1\leq k_2\leq
k_3\leq\cdots$, satisfying
\[
\begin{cases}
k_m\equiv X \bmod \varphi(p^m) & \text{if } p\neq 2 \\
k_m\equiv X \bmod 2^{m-2} & \text{if } p=2 \text{ and } m\geq 3
\end{cases}
\]
and a sequence of \uline{true} modular forms $f_i$ of weight $k_i$ and level $n$,
holomorphic at $\infty$, defined over $\W(\F_q)$, such that
\[
\begin{cases}
f_m\equiv f \bmod p^m \text{ in q-expansion}, & \text{if } p\neq 2 \\
f_m\equiv f \bmod 2^{m-1} \text{ in q-expansion} & \text{if } p=2,\ m\geq 3 .
\end{cases}
\]
Conversely. Let $\{k_m\}_{m\geq 1}$ be an \uline{arbitrary} sequence of integers,
and suppose given a sequence $f_m\in S(\W(\F_q),1,n,k_m)$ of \uline{p-adic
modular forms} of integral weights $k_i$ such that
\[
\begin{cases}
f_{m+1}\equiv f_m \bmod p^m \text{ in q-expansion at each cusp} \\
f_m\not\equiv 0 \bmod p^m \text{ in q-expansion.}
\end{cases}
\]
Then the sequence of weights $k_m$ converges to an element
$X\in\End(\Z_p^\times)$, and there is a unique modular form $f=\text{"lim"}f_m$
of weight $X$ and level $n$, holomorphic at $\infty$, defined over $\W(\F_q)$,
such that $f_m\equiv f \bmod p^m$ in q-expansion.
\end{kathm}

\begin{kathm}{Corollary}{4.5.2.}
(Serre) If a collection of elements of $\W[[q]]$ is the set of q-expansions of a
p-adic modular form $f$ of weight $X\in\End(\Z_p^\times)$ (resp. $X\in\Z_2$ if
$p=2$) and level $n$, holomorphic at $\infty$ and defined over $\W(\F_q)$, then
both $f$ and $X$ are uniquely determined.
\end{kathm}

\begin{kaproof}[Proof of the theorem]
The first part follows directly from the definitions. For the second part, we
will reduce to the case in which the $f_m$ are all \uline{true} modular forms,
whose weights satisfy $0\leq k_1\leq k_2\leq\cdots$. Indeed, if we replace $f_m$
by $f_m'=f_m(E_{p-1})^{(p^{m-1})N_m}$ with $N_m\gg 0$, then we may suppose all
$k_m\geq 0$, and, by \eqref{ka:2.7.2}, for $N_m\gg 0$, $f_m'$ has q-expansion mod $p^m$ of
a \uline{true} modular form. Rechoosing the $N_m$ to be sufficiently increasing
with $m$, we have the desired reduction. Now consider the \uline{limit}
q-expansions. We may and will work on each irreducible component of
$\Mbar_n\tens\W(\F_q)$ separately. If on a given component, the limit q-expansion
is identically zero at any cusp, it is so at every cusp, hence each $f_m$ is
$\equiv 0\ (p^m)$ on that component, and there is nothing to prove. In the
contrary case, the limit q-expansion is divisible by $p^{m_0}$ but not by
$p^{m_0+1}$ at each cusp ($m_0$ is \uline{independent} of the choice of cusp on
each irreducible component: cf.\eqref{ka:2.7.1}). Then for $m>m_0$, $f_m=p^{m_0}g_m$ where
$g_m$ is a \uline{true} modular form with q-expansions $\not\equiv 0\ (p)$. So
replacing the sequence $f_m$ by the sequence $\{f_m'\}=\{g_{m_0+m}\}$, we may
suppose that each $f_m$ has all q-expansions $\not\equiv 0 \bmod p$. Then by
\eqref{ka:4.4.1}, the congruence $f_{m+1}\equiv f_m \bmod p^m$ in q-expansion implies
that $k_{m+1}\equiv k_m$ modulo $\varphi(p^m)$ for $p\neq 2$, and modulo
$2^{m-2}$ if $p=2$ and $m\geq 3$, and that
$f_{m+1}\equiv f_m\cdot(E_{p-1})^{(k_{m+1}-k_m)/(p-1)}$ modulo $p^m$. Hence
$X=\lim k_m$ exists in $\End(\Z_p^\times)$, and
$\{f_m \bmod p^m\}_{m\geq 1}$ define a compatible family of sections of the
sheaves $\w^X$ on the schemes $S_m^\zeta$.
\end{kaproof}

\begin{kathm}{Corollary}{4.5.3.}
(Serre) Let $X\in\End(\Z_p^\times)$, and suppose $X\in\Z_2$ if $p=2$. Let
$0\leq k_1\leq k_2\leq\cdots$ be a sequence of integers such that $k_m\equiv X$
modulo $\varphi(p^m)$ if $p\neq 2$, and modulo $2^{m-2}$ if $p=2$ and $m\geq 3$.
Let $\{f_m\}$ be a sequence of \uline{true} modular forms of weight $k_m$ and
level $n$ on $\Gamma_0(p)$, holomorphic at the unramified cusps, and defined over
the fraction field $K$ of $\W(\F_q)$. Suppose that the \uline{non-constant}
terms of all the q-expansions of the $f_m$ are in $\W(\F_q)$, and that at
\uline{each} cusp,
\[
f_{m+1}(q)-f_{m+1}(0) \equiv f_m(q)-f_m(0) \bmod p^m .
\]
Then if $X\neq 0$, let $m_0$ be the largest integer such that
$X\equiv 0 \bmod \varphi(p^{m_0})$ if $p\neq 2$, for $p=2$, $m_0=1$ if $X$ is
invertible in $\Z_2$, and $m_0=2+\ord_2(X)$ if $X$ is not invertible in $\Z_2$.
Then for $m\geq m_0$, $p^{m_0}f_m$ has integral ($\in\W(\F_q)$) q-expansions, and
at each cusp we have the congruence on constant terms:
$p^{m_0}f_{m+1}(0)\equiv p^{m_0}f_m(0) \bmod p^{m-m_0}$ for all $m>m_0$ if
$p\neq 2$, and $2^{m_0}f_m(0)\equiv 2^{m_0}f_m(0) \bmod 2^{m-1-m_0}$ if $m\geq 3$
and $m\geq m_0$.
\end{kathm}

\begin{kaproof}
The integrality of the q-expansions of the $p^{m_0}f_m$ follows from \eqref{ka:4.4.3}.
Let $g_m=p^{m_0}f_m$, which has integral q-expansions. Then $g_m$ and
$h_m\xlongequal{\text{defn}}g_m\cdot(E_{p-1})^{(k_{m+1}-k_m)/(p-1)}$ have
q-expansions which are congruent modulo $p^m$ if $p\neq 2$, (resp. modulo
$2^{m-1}$ if $p=2$) and $g_m(0)=h_m(0)$. Thus $g_{m+1}-h_m$ has q-expansions
congruent to the constants $g_{m+1}(0)-h_m(0)$ modulo $p^m$ if $p\neq 2$,
(resp. modulo $2^{m-1}$ if $p=2$). Applying \eqref{ka:4.4.3} to the function
$(g_{m+1}-h_m)/p^m$ for $p\neq 2$, (resp. to $g_{m+1}-h_m/2^{m-1}$ for $p=2$) we
find that its constant term has \uline{denominator} at most $p^{m_0}$. Thus
$g_{m+1}(0)\equiv h_m(0)$ modulo $p^{m-m_0}$ if $p\neq 2$, and $2^{m-1-m_0}$ if
$p=2$.
\end{kaproof}

\begin{kathm}{Example}{4.5.4.}
(Serre) Take $f_m=G_{k_m}$, the classical Eisenstein series of level 1, whose
q-expansions are given by $-(b_{k_m})/2k_m+\sum_{n\geq 1}\sigma_{k_m-1}(n)q^n$.
Choose the $k_m$ to be \uline{strictly increasing} with $m$, so that they tend
archimedeanly to $\infty$. One checks immediately that the hypotheses of \eqref{ka:4.5.3}
are verified. The limit $\text{"lim"}p^{m_0}f_m\xlongequal{\text{defn}}
p^{m_0}G_X^*$ is thus a modular form of weight $X=\lim k_m$, whose q-expansion
is given by
\[
G_X^*(q) = \mathcal{L}^*(X) + \sum_{n\geq 1}q^n\sum_{d\mid n,p\nmid d}d^{X}(d)/d
\]
where $\mathcal{L}^*(X)$ is the (prime to $p$ part of the) Kubota-Leopoldt zeta
function, in the notation of Serre \cite{serre-proceedings}. We hasten to point out that even if the
character $X$ is an even \uline{positive integer} $2k\geq 4$, the above defined
$G_{2k}^*$ is a p-adic modular form of weight $2k$, but it is \uline{not} the
usual Eisenstein series $G_{2k}$. Indeed, the q-expansion of $G_{2k}^*$ is given
by
\[
G_{2k}^*(q) = \tfrac{1}{2}(1-p^{2k-1})\zeta(1-2k) + \sum_{n\geq 1}q^n\sum_{d\mid n,p\nmid d}d^{2k-1}
\]
while the q-expansion of $G_{2k}$ is given by
\[
G_{2k}(q) = \tfrac{1}{2}\zeta(1-2k) + \sum_{n\geq 1}q^n\sum_{d\mid n}d^{2k-1} .
\]
\uline{Both} $G_{2k}$ and $G_{2k}^*$ are p-adic modular forms of weight $2k$,
which, as Serre explained to me, are related as follows:
\[
\begin{cases}
G_{2k}^* = G_{2k} - p^{2k-1}\varphi(G_{2k}) \\
G_{2k} = \sum_{n\geq 0}(p^{2k-1}\cdot\varphi)^n(G_{2k}^*) .
\end{cases}
\]
Taking $k=1$, we obtain a p-adic modular form $G_2^*$ of weight 2, and we may
\uline{define} $G_2$ as a p-adic modular form by setting
\[
G_2 \xlongequal{\text{defn}} \sum_{n\geq 0}p^n\varphi^n(G_2^*) .
\]
An immediate calculation gives the q-expansion of $G_2$ (cf. \ref{ka:A1.3} for the series
$P$)
\[
G_2(q) = \frac{-1}{24} + \sum_{n\geq 1}q^n\sum_{d\mid n}d = \frac{-1}{24}P(q)
\]
and shows that, for any prime $p$, the series $P(q)$ is the q-expansion of a
p-adic modular form of weight two and level one. We refer the reader to \ref{ka:A2.4} for
an "intrinsic" proof of this fact for $p\neq 2,3$, based on the classical
interpretation of $P$ as a ratio of periods (cf.\ref{ka:A1.3}).
\end{kathm}
